% Generated draft: requires mathematical and source review.
\documentclass[11pt]{article}
\usepackage[margin=25mm]{geometry}
\usepackage{fontspec}
\setmainfont{Latin Modern Roman}
\usepackage{amsmath,amssymb,mathtools,unicode-math}
\setmathfont{Latin Modern Math}
\usepackage{enumitem,needspace,xurl,hyperref}
\setlength{\parindent}{0pt}
\setlength{\parskip}{5pt}
\setlength{\emergencystretch}{3em}
\newcommand{\N}{\mathbb{N}}
\newcommand{\Z}{\mathbb{Z}}
\newcommand{\Q}{\mathbb{Q}}
\newcommand{\R}{\mathbb{R}}
\newcommand{\C}{\mathbb{C}}
\newcommand{\F}{\mathbb{F}}
\newcommand{\A}{\mathbb{A}}
\newcommand{\PP}{\mathbb{P}}
\newcommand{\E}{\mathbb{E}}
\newcommand{\eps}{\varepsilon}
\newcommand{\dd}{\,\mathrm{d}}
\begin{document}
\section*{UnsolvedMath ID: 30000071}
\textbf{Source problem code:} catalogue entry
\subsection{Definitions and mathematical statement}
The source asks which hypotheses on an interaction potential force energy minimizers to crystallize. It does not fix the dimension, interaction class, particle constraints, thermodynamic regime, meaning of crystal, or quantifiers over minimizers, so it describes a program rather than one specified theorem.

A standard pair-potential instance fixes \(d\geq1\), a radial potential \(v:(0,\infty)\to\R\cup\{+\infty\}\), and distinct particle positions \(x_1,\ldots,x_N\in\R^d\), with
\[
 E_N(x_1,\ldots,x_N)=\sum_{1\leq i<j\leq N}v(|x_i-x_j|).
\]
This formula is one explicit framework, not an assumption that the original excludes anisotropic or many-body interactions. Confinement, boundary conditions, prescribed density, and hypotheses ensuring a finite energy infimum or existence of minimizers must be selected for any precise result.

An infinite periodic crystal may be modeled by \(\Lambda+B\), where \(\Lambda\) is a full-rank lattice in \(\R^d\) and \(B\) is a finite basis; a Bravais lattice is the special one-point basis case. Nonempty finite clusters cannot themselves be infinite periodic sets. Accordingly, different crystallization conclusions include exact finite-cluster containment in a crystal, local convergence of bulk configurations to a periodic set, or an asymptotically negligible fraction of defects. The source does not choose one of these inequivalent meanings, nor whether every minimizer, some minimizer, or some sequence/subsequence should have the property.

The requested task is to characterize potential hypotheses yielding the chosen crystallization conclusion in the chosen many-particle regime. No particular potential or universal crystalline lattice is proposed by the record. The supplied triage dated 21 August 2026 reports rigorous results for important two-dimensional models and an unresolved corresponding three-dimensional pair-potential problem. Its exact potential classes and theorem hypotheses are not reproduced, so it is not treated as a claim that all two-dimensional potentials crystallize or that no three-dimensional model has a crystallization result. The broad extracted question is retained as unresolved under that supplied assessment.
\subsection{Short English statement}
The source seeks conditions making many-particle energy minimizers crystalline. A precise version must choose the interaction, dimension, large-particle regime, and crystalline conclusion. The supplied curation reports specific two-dimensional progress without a general three-dimensional pair-potential resolution.
\subsection{Sources}
\begin{itemize}
\item Local UnsolvedMath registry, permanent ID \(30000071\), exact source code \texttt{catalogue entry}, Oberwolfach Reports Open Problems. Statement status \texttt{corrected\_verified}; the supplied verification reports a self-contained wording edit against the official report.
\item ``Phasen\"uberg\"ange (2005),'' Oberwolfach Reports, pp. 1587--1656, DOI \texttt{10.4171/owr/2004/30}: \url{https://doi.org/10.4171/owr/2004/30}. Proposed year and DOI refer to 2004, while the citation label says 2005.
\item Supplied ``two-dimensional crystallization theorem'': \url{https://doi.org/10.1007/s00220-005-1458-7}.
\item Supplied ``crystallization conjecture review'': \url{https://arxiv.org/abs/1504.01153}. The other supplied source is the original report DOI above; no full author/theorem details are included.
\item Literature checked 21 August 2026; supplied status \texttt{open}. Registry created 25 August 2026 and updated 16 September 2026. This dated curation triage requires independent verification before citation.
\item Registry curation attribution: Ulam AI / UnsolvedMath Contributors, CC BY 4.0. Underlying report and literature retain their supplied provenance; this statement-only draft requires human review.
\end{itemize}
\end{document}
